\title{Byte Latent Transformer: Patches Scale Better Than Tokens}

\begin{abstract}
We present the Byte Latent Transformer ({\textsc BLT}), a novel byte-level large language model (LLM) architecture that, for the first time, achieves performance on par with token-based LLMs at scale while delivering notable gains in inference speed and robustness. {\textsc BLT} transforms input bytes into patches of variable length, using these patches as the fundamental computational elements. The patch boundaries are determined according to the entropy of the subsequent byte, directing additional computational resources and model capacity to regions of greater data complexity. We report the first scaling analysis of byte-level models that controls for floating point operations (FLOPs), extending up to 8B parameters and 4T training bytes. Our findings affirm that training models directly on unprocessed bytes—without relying on a fixed vocabulary—can scale effectively. Both the training process and inference become more efficient by adaptively using longer patches in segments with higher predictability, leading to qualitative gains in reasoning and improved generalization on rare cases. In summary, {\textsc BLT} achieves significantly superior scaling compared to tokenization-based models under fixed inference costs, through concurrent expansion of both patch size and model size.
\end{abstract}

\section{Introduction}
\label{section:intro}

We present the Byte Latent Transformer~(\textbf{BLT}), a novel model that dispenses with tokenization entirely by directly learning from unprocessed byte sequences. Remarkably, BLT is the first architecture of its kind to achieve parity with traditional token-based models at scale, while offering substantial gains in both computational efficiency and robustness (\S\ref{sec:robustness}). 

Although most large language models (LLMs) are trained end-to-end, tokenization remains an exception, functioning as a heuristic preprocessing procedure that segments byte streams into a predetermined vocabulary of tokens. This imposed segmentation can introduce encoding biases that affect the model's data compression strategy and result in several limitations, including sensitivity to specific domains or modalities~\citep{dagangetting}, heightened vulnerability to input perturbations~(\S\ref{sec:robustness}), insufficient orthographic generalization~\citep{edman2024cute}, and disparities in multilingual processing~\citep{liang2023xlm,petrov2024language,limisiewicz2024myte}.

In the past, tokenization has been a crucial step because attempting to train llms directly on raw bytes incurs significant computational expense at scale, largely due to the resulting long input sequences~\citep{xue2022byt5}. Previous research has attempted to address this by introducing self-attention mechanisms that are computationally more efficient~\citep{el2020characterbert, clark2022canine} or by developing architectures that eliminate attention altogether~\citep{wang2024mambabyte}~(\S\ref{section:related_work}). Nonetheless, these strategies predominantly benefit the training of \textit{small models}. For larger models, the main source of computational burden within the Transformer architecture lies in the extensive feed-forward network layers, which process every input byte, whereas the attention mechanism contributes far less to overall costs.

To optimize compute distribution, we introduce a dynamic and trainable approach that clusters bytes into \textit{patches}~(\S\ref{section:patching}), alongside a novel architecture designed to integrate information from both bytes and patches. In contrast to tokenization, {\textsc BLT} does not rely on a static patch vocabulary. Instead, flexible byte groupings are transformed into latent patch embeddings through compact learned encoder and decoder components. Our results demonstrate that this strategy allows for a \textit{greater} efficiency in compute utilization compared to models dependent on tokenization.

Tokenization-based LLMs assign an equal amount of computational effort to each token. This approach prioritizes uniformity over efficiency, as the underlying tokenization schemes use compression strategies that do not always align with the true difficulty of predicting each token. Our model is founded on the principle that computation should be adaptively distributed, focusing resources where prediction is most challenging. For instance, the latter parts of words generally require less modeling capacity, as predicting them is typically straightforward and less uncertain compared to selecting the initial word of a sentence. This principle informs the design of {\textsc BLT}} (\S\ref{section:architecture}), which utilizes three transformer modules: two lightweight, byte-level \textit{local models} and one large-scale, global \textit{latent transformer} (\autoref{fig:arch_high_level}). 

{\textsc BLT} dynamically forms byte clusters by segmenting the input according to the entropy associated with next-byte predictions. This strategy yields contextual groupings with more consistent levels of informational complexity, allowing the model to more effectively allocate computation.

We conduct the first study examining the scaling behavior of byte-level models under a fixed flop budget, training models as large as 8B parameters on 4T bytes of data without relying on fixed-vocabulary tokenization for end-to-end training. Our results demonstrate that {\textsc BLT} achieves comparable performance to Llama 3 under matched training computational budgets,\footnote{The computational requirement is assessed by tallying the total number of Floating Point OPerations (flops) involved.} while offering up to a 50\% reduction in flops during inference~(\S\ref{section:scaling}). Furthermore, we find that processing inputs at the byte level leads to notable improvements in capturing rare or long-tail linguistic phenomena. {\textsc BLT} exhibits increased robustness to input noise and stronger performance on character-level processing tasks, as shown in domains such as orthographic recognition, phonological analysis, and low-resource machine translation~(\S\ref{sec:robustness}). Moreover, {\textsc BLT} enables concurrent scaling of model size and patch length without exceeding the inference flop constraints; utilizing longer patch sizes typically lowers computational costs, which can then be allocated towards expanding the global latent transformer component, as it is invoked less frequently. In inference-flop-constrained scaling experiments~(\autoref{fig:fixed_inference_scaling}), we find that this approach produces substantially improved scaling patterns relative to traditional tokenization-based models.

To conclude, our work offers the following key contributions: 1) We present {\textsc BLT}, a byte latent LLM framework that adaptively assigns computational resources to enhance FLOP efficiency, 2) We demonstrate that our method attains training FLOP-controlled equivalence with Llama 3 up to the 8B parameter range, with the flexibility to trade slight reductions in evaluation performance for up to 50\% greater FLOP efficiency, 3) The {\textsc BLT} approach opens up a novel path for LLM scaling, enabling increases in model size while holding inference costs constant, and 4) We provide evidence that {\textsc BLT} models show heightened resilience to input noise and possess sensitivity to sub-word features in input data that are typically overlooked by token-centric LLMs. The training and inference code for {\textsc BLT} is publicly available at~\url{https://github.com/facebookresearch/blt}.

\section{Patching: From Individual Bytes to Groups of Bytes}
\label{section:patching}

Dividing byte sequences into \textit{patches} enables {\textsc BLT} to flexibly assign computational resources according to the contextual requirements. Figure~\ref{fig:patching_types} illustrates multiple approaches for organizing bytes into patches. More formally, a patching function $f_p$ partitions a byte sequence $\pmb{x}=\{x_i,|i=1,\ldots n\}$ of length $n$ into $m < n$ patches $\pmb{p}=\{p_j|j=1,\ldots,m\}$ by designating each $x_i$ as belonging to the set \{0,1\}, with 1 marking the onset of a new patch. In both token-based and patch-based architectures, the key determinant of computational expense is the number of main Transformer operations, which, in the context of {\textsc BLT}, corresponds to the quantity of patches the patching function yields for the data. As a result, the mean size of a patch—referred to as the \textit{patch size}—becomes the primary metric for estimating compute costs associated with a specific patching strategy during both model training and inference~(\S\ref{section:flops}). 

We next describe three specific patching schemes: patching using a constant byte count per patch~(\S\ref{section:static-patch}), dividing at whitespace characters~(\S\ref{section:white-patch}), and adaptive patching that leverages entropy information from a lightweight byte language model~(\S\ref{section:dyn-patch}). Finally, we consider incremental patching methods and explore the distinctions between patching and tokenization~(\S\ref{section:bpe-patch}).

\subsection{Strided Patching Every K Bytes}
\label{section:static-patch}

One of the simplest methods for grouping bytes is to partition them into fixed-size patches of length $k$, as implemented in MegaByte~\citep{yu2023megabyte}. This approach, utilizing a uniform stride, is straightforward to realize in both training and inference pipelines, and facilitates easy adjustment of the average patch size, thereby offering direct control over computational costs. Nonetheless, several important limitations are inherent to this patching strategy. Firstly, computational resources are not adaptively distributed according to informational content: for instance, a transformer step $j$ might be underutilized when processing uninformative segments like whitespace in code, or, conversely, might not provide enough capacity for segments rich in content, such as mathematical expressions. Secondly, this fixed approach results in inconsistent, non-context-aware segmentation of similar byte sequences, causing, for example, identical words to be divided at arbitrary points.

\subsection{Space Patching}
\label{section:white-patch}

\citet{slagle2024spacebyte} introduces an enhancement to strided patching by generating new patches whenever a space-like byte\footnote{
    Space-like bytes refer to any byte that does not correspond to a Latin letter, a digit, or a UTF-8 continuation byte. Furthermore, each patch must contain at least one byte that is not space-like. } appears, as these typically serve as boundaries for linguistic units in various languages. In the Space patching approach, each word receives dedicated computation through a latent transformer step (i.e., increased FLOPs), thereby standardizing the patching of words across different sequences and allocating greater computational resources to challenging predictions that usually arise after spaces. For instance, guessing the initial byte of a response to the prompt ``Who composed the Magic Flute?\uline{\hspace{.6em}}'' is notably more challenging than continuing the completion after ``M'', because the set of probable options narrows significantly once the first character is determined—making the rest of the answer, such as ``Mozart'', easier to predict. Despite these advantages, space patching does not effectively address all languages and domains and, crucially, lacks flexibility in patch size. To address these issues, we present an alternative patching approach that builds on the observation that word-initial bytes are often the most challenging to predict, while simultaneously affording a straightforward means of adjusting patch size.

\subsection{Entropy Patching: Using Next-Byte Entropies from a Small Byte LM}
\label{section:dyn-patch}

Instead of utilizing a heuristic based on fixed rules like whitespace, our method adopts a data-driven strategy to detect regions where the model exhibits high uncertainty in predicting the next byte. We propose \textit{entropy patching}, a technique that leverages entropy measurements to determine the limits of patches.

To implement this, we fit a compact byte-level auto-regressive language model to the {\textsc BLT} training corpus, then calculate the entropy of the next byte predictions given the LM-induced distribution $p_{e}$ over the set of possible bytes $\mathcal{V}$:

\begin{equation}
H(x_i) = \sum_{v \in \mathcal{V}} p_{e}(x_i=v|\pmb{x}_{<i}) \log p_{e}(x_i=v|\pmb{x}_{<i})
\end{equation}

We employ two strategies to detect patch boundaries based on entropies $H(x_i)$. The initial method selects points that exceed a universal entropy threshold, as depicted in~\autoref{fig:patching}. The alternative technique highlights locations where the entropy is significantly higher compared to the immediately preceding value. This latter strategy can also be understood as pinpointing instances where the generally decreasing entropy trend within a patch is disrupted.

\begin{align*}
    \text{Global Constraint}\hspace{1cm}H(x_t)& > \theta_g\\
    \text{Approx. Monotonic Constraint}\hspace{1cm}H(x_t)& - H(x_{t-1}) > \theta_r
\end{align*}

During dataloader initialization, patch boundaries are detected through an efficient preprocessing procedure. This approach contrasts with that of~\citet{nawrot-etal-2023-efficient}, who employ a classifier trained to determine patch boundaries based on entropy estimates. In our experiments~(\S\ref{section:experiments}), we evaluate and compare both strategies for separating bytes with low versus high entropy.

\subsection{The Byte-Pair Encoding (BPE) Tokenizer and Incremental Patching}
\label{section:incremental-patching}
\label{section:bpe-patch}

Numerous contemporary llms, such as our reference model Llama 3, implement subword tokenization methods like bpe~\citep{gage1994new, sennrich-etal-2016-neural}. In our discussion, ``tokens'' denote groups of bytes selected from a \textit{finite} vocabulary that is specified before training, while ``patches'' indicate adaptively defined byte sequences that lack a predetermined vocabulary. An important distinction between tokens and patches lies in the fact that when working with tokens, the model does not have immediate access to the original byte-level information.

An important advancement introduced by {\textsc BLT} compared to traditional tokenization-based architectures is its ability to reconceptualize the balance between vocabulary magnitude and computational requirements. In conventional large language models, expanding the vocabulary typically results in larger average token sizes, which consequently reduces the number of inference steps needed. However, this also necessitates a broader output dimension in the model’s terminal projection layer. As a result, tokenization-driven methods face significant constraints in varying token size without disproportionately impacting inference efficiency and resource demands. For instance, Llama 3 achieves a rise in average token size from 3.7 to 4.4 bytes, but this improvement comes at the expense of quadrupling the embedding table’s size relative to Llama 2.

During generation, {\textsc BLT} must determine at each position in the byte sequence whether a patch boundary has been reached, as this dictates if additional computation by the Latent Transformer is necessary. Importantly, this determination must be made solely based on the already-generated portion of the sequence, without reliance on as-yet-unproduced future bytes. As a result, the method of patching must not presuppose knowledge of forthcoming bytes for making segmentation decisions. More precisely, a patching strategy $f_p$ adheres to the principle of incremental patching if the following holds:
$$f_p(\pmb{x}_{<i}) = f_p(\pmb{x})_{<i}$$
bpe fails to constitute an incremental patching strategy, since an identical prefix could be segmented differently contingent upon the succeeding bytes, thus violating the criterion stated above\footnote{Introducing a distinct delimiter token to demarcate patch boundaries can render bpe incremental, but at the cost of elongating the byte sequence.}.

\section{{\textsc BLT} Architecture}
\label{section:architecture}
{\textsc BLT} integrates a large-scale, global autoregressive language model tasked with processing patch-level representations. This is complemented by two smaller, local modules responsible for converting byte sequences into patches and reconstructing bytes from patch representations, as depicted in~\autoref{fig:arch_high_level}.

\subsection{Latent Global Transformer Model}

\textit{The Latent Global Transformer} functions as an autoregressive transformer, denoted by $\mathcal{G}$ and comprised of $l_{\mathcal{G}}$ layers, which transforms a series of latent input patch representations $p_j$ into corresponding output patch representations $o_j$. Within this work, the subscript $j$ indicates patches while $i$ refers to bytes. For attention, the global model employs a block-causal attention mask~\citep{dubey2024llama}, enforcing that each patch attends only to itself and preceding patches in the current document. The computational burden during both pre-training and inference is largely due to this model, making the decision of when to utilize it a means of dynamically allocating computational resources based on the complexity of the input and output regions.

\subsection{Local Encoder}
\label{section:local}

\textit{The Local Encoder Model}, represented as $\mathcal{E}$, utilizes a compact transformer framework characterized by significantly fewer layers, $l_{\mathcal{E}} << l_{\mathcal{G}}$. Its principal function is to swiftly transform input byte sequences $b_i$ into rich patch-level representations, $p_j$. Distinctively, this architecture integrates a cross-attention mechanism after each transformer layer, designed to aggregate byte-level features into patch representations~(\autoref{fig:crossattn}).

Initially, each byte $b_i$ from the input is mapped to an embedding $x_i$ through a lookup in a parameter matrix of dimension $\mathbb{R}^{256 \times h_{\mathcal{E}}}$. These embeddings may optionally be enriched with hash-embedding information~(\S\ref{sec:hash-emb}). Subsequently, the model applies a stack of alternating transformer layers and cross-attention layers~(\S\ref{sec:enc-cross-attn}), iteratively refining the representations to produce the patch embeddings $p_i$, which are then provided to the global transformer module $\mathcal{G}$.

The transformer portion adopts a \textit{local block causal} attention mask: each byte is restricted to attend within a fixed window of $w_{\mathcal{E}}$ previous bytes, permitting attention across dynamic patch boundaries but prohibiting access across distinct document boundaries. Subsequent sections elaborate further on the embedding approach and the architecture of the cross-attention layer.

\subsubsection{Encoder Hash n-gram Embeddings}
\label{sec:ngram-emb}
\label{sec:hash-emb}
Capturing information from previous bytes is essential for building strong and expressive representations at every position $i$. In {\textsc BLT}, this is accomplished by representing not only the current byte $b_i$ itself, but also its context within a byte n-gram. Specifically, for each position $i$, we begin by forming byte-grams

\begin{align}
    g_{i,n}=\{b_{i-n + 1},\ldots, b_i\}
\end{align}

for every byte position $i$ and for $n$ ranging from three to eight.\footnote{
Byte-grams of length $n$ or greater are excluded in cases where $i<n$. }

Subsequently, we propose the use of hash $n$-gram embeddings, where each byte $n$-gram is mapped through a hash function to an index within a dedicated embedding table $E_{n}^{hash}$, with fixed dimensionality, corresponding to each $n \in \{3,4,5,6,7,8\}$~\citep{bai2010learning}.
The generated embedding is incorporated into the byte’s own embedding, after which the sum is normalized and supplied as input to the local encoder model. The final, enhanced embedding is computed as follows:

\begin{align}
e_i &= x_i + \sum_{n = 3,...,8}E_{n}^{hash}(\text{Hash}(g_{i,n})) \\
\text{where, Hash}(g_{i,n}) &= \text{RollPolyHash}(g_{i,n}) \% |E_{n}^{hash}|
\end{align}

Each $e_i$ is normalized by dividing by the total count of $n$-gram sizes incremented by one, utilizing the RollPolyHash algorithm as detailed in Appendix~\ref{appendix:rollpolyhash}. Section~\ref{section:ablations} examines the impact of varying both $n$ and the size of the embedding table in $n$-gram hash embeddings on scaling laws under flop constraints. Beyond hash-based $n$-gram embeddings, we also investigated the use of frequency-based $n$-gram embeddings; further information regarding these experiments can be found in Appendix~\ref{appendix:ngrams}.

\subsubsection{Encoder Multi-Headed Cross-Attention}
\label{sec:enc-cross-attn}
Our approach largely mirrors the cross-attention mechanism employed in the input stage of the Perceiver model~\citep{jaegle2021perceiver}; however, a key distinction is that our latent representations are derived from the representations of variable-sized patches rather than relying on a predetermined set of latent variables~(\autoref{fig:crossattn}). Each latent representation focuses exclusively on the bytes within its associated patch. For each patch $p_j$, a corresponding query vector is constructed by aggregating the byte-level embeddings specific to $p_j$. This aggregated representation is then linearly projected via $\mathcal{E}_{C} \in \mathbb{R}^{h_{\mathcal{E}} \times (h_{\mathcal{E}} \times U_{\mathcal{E}})}$, where $U_{\mathcal{E}}$ represents the number of heads in the encoder cross-attention layer. Denoting by $f_{\text{bytes}}(p_j)$ the ordered set of bytes within the patch $p_j$, we then compute

\begin{align}
P_{0,j} &= \mathcal{E}_{C}(f_\text{bytes}((p_j)), f~ \text{is a pooling function} \\
P_l &= P_{l-1} + W_o\left(\text{softmax}\left(\frac{QK^T}{\sqrt{d_k}}\right)V\right) \\
\text{where } Q_j &= W_q(P_{l-1,j}), K_i = W_k(h_{l-1,i}), V_i = W_v(h_{l-1,i}) \\
h_l &= \text{Encoder-Transformer-Layer}_l(h_{l-1})
\end{align}

Here, $P \in \mathbb{R}^{n_p \times h_{\mathcal{G}}}$ denotes the $n_p$ patch vectors intended for input to the global model. The initialization of these representations is performed by aggregating the byte embeddings $e_i$ that make up each individual patch $p_j$. The matrices $W_q$, $W_k$, $W_v$, and $W_o$ are used as the respective linear projections for queries, keys, values, and outputs. The keys and values specifically derive from projections of the byte-level representations $h_i$ originating from the preceding layer (in the first layer, these are given by $e_i$). A tailored masking approach is adopted for patching, restricting each query $Q_j$ to interact only with the keys and values linked to the bytes within its own patch $j$. As multi-headed attention is applied to $Q$, $K$, and $V$, and given that the patch representation dimensionality ($h_{\mathcal{G}}$) generally exceeds $h_{\mathcal{E}}$, we preserve $P_l$ as several attention heads each of size $h_{\mathcal{E}}$ for cross-attention. These are subsequently concatenated to yield vectors of dimension $h_{\mathcal{G}}$. Furthermore, queries, keys, and values undergo pre-LayerNorm normalization, and no positional encodings are introduced in this cross-attention component. The output of the cross-attention module also incorporates a residual connection.

\subsection{Local Decoder} 

In line with the local encoder, the local decoder $\mathcal{D}$ is implemented as an efficient transformer architecture, consisting of $l_{\mathcal{D}}$ layers where $l_{\mathcal{D}} << l_{\mathcal{G}}$. Its function is to reconstruct raw bytes, $y_i$, from a series of global patch representations, $o_j$. By conditioning on previously generated bytes, the local decoder generates each element of the output byte sequence, relying on hidden states supplied by the local encoder for the corresponding byte inputs. The decoding process involves $l_{\mathcal{D}}$ layers arranged in an alternating pattern of cross-attention followed by transformer operations. Within each layer, cross-attention is utilized prior to the transformer block, enabling the construction of byte-level representations from patch-level information, after which the transformer component processes these byte representations.

\subsubsection{Decoder Multi-headed Cross-Attention} 

Within the decoder’s cross-attention mechanism, the assignment of queries and key/values is reversed: byte-representations serve as queries, while patch representations function as the key/values. For the cross-attention operation, the byte-representations are initially set using the byte embeddings output by the final encoder layer, specifically $h_{l_{\mathcal{E}}}$. For each subsequent layer $l$, the byte-representations $d_{l,i}$ are updated according to:

\begin{align}
D_0 &= h_{l_{\mathcal{E}}} \\
B_l &= D_{l-1} + W_o\left(\text{softmax}\left(\frac{QK^T}{\sqrt{d_k}}\right)V\right), \\
  &\text{where } Q_i = W_q(d_{l-1,i}), K_i = W_k(\mathcal{D}_C(o_j)), V_i = W_v(\mathcal{D}_C(o_j)) \\
  D_l &= \text{Decoder-Transformer-layer}_l(B_l)
\end{align}

Here, $W_k$ and $W_v$ refer to the key and value projection matrices, respectively; these matrices perform linear transformations along with a split operation $\mathcal{D}_C$ on the terminal patch representations $o_j$ output by the global model. The matrix $W_q$ projects queries derived from the byte-level representations $d_{l-1}$ of the preceding decoder transformer layer (or from $h_{l_{\mathcal{E}}}$ at the initial layer). $W_o$ denotes the output projection matrix. Consequently, this yields $B \in \mathbb{R}^{h_{\mathcal{D}} \times n_b}$, with $n_b$ representing the total output bytes. The decoder transformer layer subsequently processes $B$, which is the result from the cross-attention stage, to yield the next set of decoder representations, $D_l$. Mirroring the setup in the local encoder's cross-attention, we implement multi-head attention, apply pre-LayerNorm, omit positional embeddings, and incorporate a residual connection around the cross-attention mechanism.

\section{Experimental Setup}
\label{section:experiments}
We have rigorously constructed a suite of controlled experiments to evaluate {\textsc BLT} against models based on tokenization, explicitly ensuring that {\textsc BLT} is not inadvertently favored by advantages linked to the utilization of longer input sequences.

\subsection{Pre-training Datasets} In our experiments, all model sizes are pre-trained using two main datasets: (1) The Llama 2 dataset~\citep{touvron2023llama}, consisting of 2 trillion tokens sourced from diverse, publicly accessible datasets that undergo comprehensive cleaning and filtering processes to enhance data quality; and (2) {\textsc BLT}-1T: a novel collection totaling 1 trillion tokens, curated from multiple public sources, with the addition of a portion of pre-training data shared by Datacomp-LM~\citep{li2024datacomp}. The Llama 2 dataset supports scaling law investigations, particularly following guidelines from~\cite{dubey2024llama}, to optimize the architectural configurations for {\textsc BLT}, while {\textsc BLT}-1T enables a full-scale pre-training evaluation for comparisons with Llama 3 in downstream applications. Importantly, neither dataset incorporates any information derived from Meta’s platforms or services. For experiments involving tokenizer-based models, we utilize the Llama 3 tokenizer—which features a vocabulary of 128K tokens—as it consistently yielded superior baseline results compared to the Llama 2 tokenizer within our tests.

\subsection{Entropy Model} For our experiments, the entropy model consists of a byte-level language model that is trained on the identical data distribution as the complete {\textsc BLT} framework. By default, this model utilizes a transformer architecture with 100M parameters, comprising 14 layers, a hidden size of 512, and employs sliding window attention spanning 512 bytes. All other hyperparameter settings mirror those employed in both our local and global transformer configurations. We also explored a variety of model scales, architectural designs, and receptive field sizes, as further elaborated in \autoref{section:ablations}. Notably, when the receptive field is sufficiently restricted, it becomes possible to implement the trained entropy model as an efficient lookup table.

\subsection{Entropy Threshold and Equalizing Context Length} In entropy-based patching approaches, we determine a patching threshold to control the expected average \textit{patch size} on the pretraining data mixture. Within {\textsc BLT}, as opposed to fixed tokenization, the \textit{patch size} is a flexible parameter and can be set arbitrarily, which directly affects the context size processed by the model. To ensure a fair comparison and prevent models with larger patches from benefiting from increased effective context, we hold the total number of bytes per batch constant in expectation. Consequently, models employing larger patch sizes are assigned proportionally shorter sequences. Specifically, for Llama 2 data we set an 8k byte context, whereas for the {\textsc BLT}-1T dataset, the context size is raised to 16k bytes on average, all while keeping the mean batch size fixed at 16M bytes.

Although the mean batch size does not vary, dynamic patching techniques produce varying proportions of bytes per patch when assembling data batches. To optimize efficiency, our {\textsc BLT} training system groups together batches composed of patches, which eliminates the need for padding operations in the costlier latent transformer component. Consequently, each batch contains an identical count of patches. To prevent excessive memory usage from sequences containing abnormally large patches, we pad and, if necessary, truncate byte sequences to fixed lengths of 12k and 24k bytes for Llama 2 and {\textsc BLT}-1T datasets, respectively, during training.

\subsection{Entropy Model Context}
\label{section:entropy-context}
Our empirical observations indicate that entropy patching tends to produce increasingly larger patches when applied to highly structured material, such as multiple choice questions, which are characteristically repetitive (refer to the patching example in \autoref{fig:mmlu_patching} for MMLU). These differences emerge due to the reduced entropy associated with the redundant segments present in the entropy model context. Consequently, for the large-scale {\textsc BLT}-Entropy experiment utilizing a patch size of 4.5, we refresh the entropy context using new line characters and employ an approximate monotonicity constraint, as this method demonstrates increased resilience to "entropy drift" caused by fluctuations in context length. It is important to note that this modification only influences how entropy values are computed, whereas the approach for determining the entropy threshold remains unchanged.

\subsection{FLOPs Estimation} 
\label{section:flops}

Our calculation of transformer FLOPs is primarily based on the framework outlined in Chinchilla~\citep{hoffmann2022training}, which encompasses FLOPs for the feed-forward networks, qkvo projections within the self-attention mechanism, as well as the calculation of attention scores and output projection. However, in contrast to their approach, we consider the input embedding stage as a computationally efficient lookup table rather than a dense matrix multiplication, resulting in zero FLOPs attributed to this component. Consistent with earlier studies, we assume that the backward pass entails double the FLOP count of the forward pass.

To determine the number of flops \textit{per byte} for {\textsc BLT} models, we sum the flop counts from the local encoder transformer, the global latent transformer, and the local decoder transformer, in addition to the cross attention modules present in both the encoder and decoder:

\begin{align}
    \text{FL}_{\text{{\textsc BLT}}} &= \text{Transf. FL}(h_{\mathcal{G}}, l_{\mathcal{G}}, m=n_{ctx}/n_p,V=0) / n_p\\
   &+ \text{Transf. FL}(h_{\mathcal{E}}, l_{\mathcal{E}}, m=w_{\mathcal{E}}, V=0) \\
   &+ \text{Transf. FL}(h_{\mathcal{D}}, l_{\mathcal{D}}, m=w_{\mathcal{D}}, V=256) \\ 
    &+ \text{Cross Attn. FL}(h_{\mathcal{E}}, l_{\mathcal{E}}, m=n_p, r=n_p/k) \times k/n_p \\ 
   &+ \text{Cross Attn. FL}(h_{\mathcal{D}}, l_{\mathcal{D}}, m=k, r=k/n_p) 
\end{align}

Here, $n_{ctx}$ denotes the length of the sequence in bytes, $n_p$ represents the patch size, and $r$ indicates the ratio between the number of queries and key/value pairs. The parameter $k$ defines the proportion of the patch dimension to the byte dimension, effectively representing the number of local model splits that are concatenated to build the complete model representation (with $k=2$ as shown in \autoref{fig:crossattn}). The symbol $V$ stands for the vocabulary size in the output projection layer, which is exclusively utilized within the local decoder. The attention module adopts a context length, $m$, that varies based on whether it operates on byte-level or patch-level sequences. We adjust the computation of attention FLOPs accordingly for each constituent module. Detailed formulations for calculating Transformer-FLOPs and Cross-Attention FLOPs are available in Appendix~\ref{section:flops-appendix}.

\subsection{Bits-Per-Byte Estimation} The utility of perplexity is confined to scenarios involving a consistent tokenizer, as it quantifies the ambiguity on a per-token basis. To enable direct comparison between models operating at the byte and token granularity, we adopt the Bits-Per-Byte (BPB) metric in line with prior studies~\citep{xue2022byt5, yu2023megabyte, wang2024mambabyte}, which serves as a tokenizer-agnostic counterpart to perplexity. This metric is defined as follows:

\begin{align}
    \text{BPB}(x)&= \frac{\mathcal{L}_{CE}(\pmb{x})}{\ln(2)\cdot n_{\text{bytes}}}
\end{align}

Here, the aggregate cross-entropy loss across the dataset $\pmb{x}$ is divided by both the total byte count of $\pmb{x}$ and a normalization constant, yielding a standardized measure of model uncertainty irrespective of tokenization.

\subsection{Transformer Architecture Hyperparameters} In configuring the transformer modules of {\textsc BLT}, encompassing both local and global models, we adopt design choices similar to those of Llama~3~\citep{dubey2024llama}. Specifically, the feed-forward sublayers implement the SwiGLU activation~\citep{swiglu}, while rotary positional embeddings (RoPE)~\citep{RoPe}, set with $\theta=500000$ as described in~\citep{xiong2024effective}, are incorporated exclusively within the self-attention mechanisms. Layer normalization is performed via RMSNorm~\citep{RMSNorm}. For self-attention operations employing standard attention masking—namely, \textit{block causal} and \textit{fixed-window block causal}—we utilize Flash attention~\citep{dao2022flashattention}, with the attention mask width fixed at 512. In contrast, because the cross-attention mechanisms are governed by patch-specific dynamic masks, we leverage Flex Attention\footnote{\url{https://pytorch.org/blog/flexattention}} to generate optimized kernel implementations, thereby greatly accelerating model training.

\subsection{{\textsc BLT}-Specific Hyperparameters} 

To evaluate the performance of {\textsc BLT} models, we design experiments focusing on two main aspects: scaling behavior and performance on downstream tasks. For this purpose, we include models at several parameter counts: 400M, 1B, 2B, 4B, and 8B. Detailed architectural hyperparameters corresponding to each model size are available in Appendix~Table~\ref{tab:arch}. 

For the initialization of queries in the first cross-attention module of the local encoder, we adopt a max-pooling strategy. All {\textsc BLT} models are configured with 500,000 hashes, employing a single hash function, and n-gram lengths spanning from 3 to 8. We fix the learning rate at $4e-4$ across all models. This learning rate aligns with findings from a hyperparameter sweep across values from $1e-3$ to $1e-4$ at the 400M and 1B model scales, indicating the same optimal setting for both token-based and {\textsc BLT} models. For scaling experiments utilizing the Llama-2 dataset, batch sizes are selected based on the recommendations in~\cite{dubey2024llama}, or their byte-size equivalents. 

Optimization is conducted with the AdamW algorithm~\citep{adamw}, setting $\beta_1$ to 0.9 and $\beta_2$ to 0.95, with an $\epsilon=10^{-8}$. The learning rate is scheduled with a linear warm-up over the initial 2000 steps, followed by cosine annealing down to zero. We incorporate a weight decay coefficient of 0.1 and apply global gradient clipping at a maximum norm of 1.0.

\section{Scaling Trends}
\label{section:scaling}

We provide a comprehensive overview of the scaling behaviors observed in byte-level models, with the goal of guiding subsequent scaling efforts for {\textsc BLT} models. Our scaling analysis seeks to overcome shortcomings in prior studies of byte-level models by (a) examining scaling patterns under the compute-optimal training setup, (b) training aligned 8B parameter models using substantial training corpora—up to 1T tokens or 4T bytes—and assessing performance on downstream benchmarks, and (c) analyzing scaling laws within fixed inference-cost scenarios. Further on, we will explore particular benefits that arise from modeling data as byte-sequences.

\subsection{Parameter Matched Compute Optimal Scaling Trends}
\label{section:parameter-matched-scaling}
Leveraging the Llama 2 dataset, we conduct experiments with a series of \textit{compute-optimal} bpe and {\textsc BLT} models at four distinct scales, specifically spanning from 1B to 8B parameters. For these models, we chart the relationship between the training flops expended and the resulting language modeling accuracy, utilizing a representative portion of the mixed training corpus. In the case of bpe models, training is performed under the configuration where the proportion of model parameters to training tokens aligns with the optimum established by Llama~3~\citep{dubey2024llama}. This approach to \textit{compute-optimality} is theoretically intended to maximize performance given a constrained training budget~\citep{hoffmann2022training}, thus furnishing a rigorous point of comparison for our study. Furthermore, for every bpe model configuration, we train a matching {\textsc BLT} model utilizing the same training dataset and equipping it with a Latent Transformer whose architecture and parameter count are aligned to its bpe counterpart.

As shown in~\autoref{fig:scaling-compute-opt} (right), {\textsc BLT} models consistently equal or surpass the performance of their bpe equivalents, with this pattern persisting as both model size and compute scale up. To our knowledge, {\textsc BLT} represents the first Transformer architecture operating at the byte level to exhibit scaling behavior comparable to that of BPE-based models in compute-efficient settings. These results therefore support our hypothesis that the parameter-to-training-compute ratio considered optimal for bpe models is also applicable to {\textsc BLT}, or at the very least, does not deviate significantly.

Advancements in model architecture, along with dynamic patching techniques, are essential for achieving competitive bpe scaling behavior. In ~\autoref{fig:scaling-compute-opt} (left), we present a comparison between models employing space-patching strategies and Llama 3. The representation of SpaceByte \citep{slagle2024spacebyte} is approximated using {\textsc BLT} space-patching, excluding both n-gram embeddings and cross-attention mechanisms. While SpaceByte demonstrates gains over Megabyte, its performance still significantly lags behind Llama 3. Meanwhile, \autoref{fig:scaling-compute-opt} (right) highlights the gains resulting from architectural enhancements in combination with dynamic patching. {\textsc BLT} models, when appropriately scaled, exhibit performance comparable to state-of-the-art models that utilize tokenizers, such as Llama 3.

Additionally, we find that tokenizer selection significantly impacts the performance of tokenizer-based models; specifically, those trained with the Llama-3 tokenizer surpass models utilizing the Llama-2 tokenizer, given identical training datasets.

Ultimately, our {\textsc BLT} model exhibits intermediate behavior between Llama 2 and 3, particularly when employing much larger patch sizes. Whereas the average token length of the Llama 2 and 3 bpe tokenizers is 3.7 and 4.4 bytes respectively, {\textsc BLT} manages to maintain similar scaling patterns with mean patch sizes of 6 or even 8 bytes. Because inference flops decrease as the average patch size increases, setting the patch size to 8 bytes could yield nearly 50\% reduction in inference flops. Moreover, models configured with larger patch sizes demonstrate enhanced performance when both the model and data scale up. For instance, {\textsc BLT} at an 8-byte patch size initially underperforms compared to bpe Llama 2 at 1B parameters, but surpasses the bpe model by the time the scale reaches 7B. This trend implies that larger patch sizes might continue to improve at still greater scales, and suggests that using even bigger patches could become advantageous as both model sizes and computational budgets expand.

\subsection{Beyond Compute Optimal Task Evaluations}
\label{section:evals}
To further investigate scaling dynamics, we train an 8B {\textsc BLT} model past the compute optimal point, utilizing the {\textsc BLT}-1T dataset—a larger, higher-quality corpus—and evaluate its capabilities across a range of established classification and generation benchmarks. The evaluation encompasses standard tasks involving common sense reasoning, general world knowledge, and code synthesis:
\paragraph{Classification tasks} consist of ARC-Easy (0-shot)~\citep{Eval_ARC}, Arc-Challenge (0-shot)~\citep{Eval_ARC}, HellaSwag (0-shot)~\citep{Eval_hellaswag}, PIQA (0-shot)~\citep{Eval_piqa}, and MMLU (5-shot)~\citep{hendrycks2020measuring}. For these benchmarks, we utilize a prompt likelihood approach, measuring the probability assigned to possible answer choices, and compute mean accuracy scores.
\paragraph{Coding related generation tasks:} To examine code generation proficiency, we present pass@1 results for MBPP (3-shot)~\citep{Eval_MBPP} and HumanEval (0-shot)~\citep{Eval_HumanEval}, gauging the model’s performance in Python code generation scenarios.

\autoref{tab:evals} presents a comparison among three models, all trained on the {\textsc BLT}-1T dataset: a bpe-based model using the Llama 3 tokenizer,\footnote{We select the Llama 3 tokenizer, featuring a 128k vocabulary, due to its superior performance relative to the 32k vocabulary of Llama 2.} and two distinct versions of the {\textsc BLT} model—one implementing a space-patching strategy ({\textsc BLT}-Space), and the other adopting an entropy-driven patching approach ({\textsc BLT}-Entropy). The latter incorporates an approximate monotonicity constraint and reinitializes its entropy context at newlines, as detailed in~\autoref{section:entropy-context}. All models are trained using roughly equivalent computational resources (flop budget). Notably, for {\textsc BLT}-Entropy, we lower the entropy threshold during inference from 0.6 to 0.1, a modification that yields improved performance across tasks, albeit requiring more inference steps.

The {\textsc BLT}-Entropy model achieves superior performance compared to the Llama 3 model on 4 of the 7 evaluated tasks, despite both models being trained with an equivalent data volume in bytes. This performance gain is likely attributable to two key factors: (1) enhanced utilization of training computation made possible by dynamic patching, and (2) the explicit modeling of information at the byte level rather than the token level.

Conversely, {\textsc BLT}-Space yields inferior results compared to the Llama 3 tokenizer across all tasks except one. Nevertheless, it offers a notable decrease in inference flops owing to its greater average patch size of 6 bytes. By contrast, both the bpe and entropy-patching models demonstrate similar average patch sizes of about 4.5 bytes when evaluated on the training data mix. Given an equivalent training budget, the model utilizing the larger patch size is able to process 30\% more data relative to the other two approaches, potentially moving {\textsc BLT} further from the compute-optimal regime.

\subsection{Patches Scale Better Than Tokens} Using {\textsc BLT} models enables concurrent scaling of both model and patch sizes without altering the training and inference flop requirements or increasing the amount of training data. Unlike fixed-vocabulary token-based models, patch-based approaches uniquely allow for unrestricted increases in patch size, thus circumventing inherent efficiency tradeoffs, as elaborated in Section~\ref{section:bpe-patch}. By employing longer patches, computational demands decrease, which in turn allows resources to be redirected toward expanding the global latent transformer, since it needs to be executed less frequently.

We perform a fixed inference scaling experiment to investigate whether increasing model size, coupled with fewer computational steps on larger patches, yields superior performance compared to smaller models requiring more steps. Using initial model scales of 400m and 3.6B parameters with the Llama 2 tokenizer, we identify models of comparable computational cost with the Llama 3 tokenizer, as well as {\textsc BLT}-Entropy architectures featuring mean patch sizes of 6 and 8 bytes within the training datamix (refer to~\autoref{tab:fixed-inf-params} for detailed model configurations). For models utilizing a patch size of 8, we implement 3 encoder layers in place of 1. Each model is trained under multiple total training flop constraints.

\autoref{fig:fixed_inference_scaling} demonstrates that {\textsc BLT} models exhibit superior scaling behavior compared to architectures reliant on tokenization across both categories of inference computational requirements. For limited training budgets, BPE-based models initially show stronger performance, but {\textsc BLT} quickly overtakes them just beyond the compute-optimal threshold. This indicates that, in practical scenarios, allocating greater resources during the one-time pretraining phase may be advantageous, yielding higher performing models under a fixed inference cost. The 8B model class, exemplified by Llama 3.1, illustrates this approach; such models are trained on data volumes that exceed by two orders of magnitude the compute-optimal amount typically recommended for their size.

The threshold at which {\textsc BLT} surpasses token-based approaches has moved slightly closer to the compute-optimal threshold when transitioning to models within a higher FLOP category—shifting from requiring 3x to only 2.5x the compute-optimal budget. In the same vein, for the model with patch size 8, we observe a steeper scaling curve in this larger FLOP regime, allowing it to surpass competing models at an earlier stage. As outlined in Section~\ref{section:parameter-matched-scaling}, increasing the patch size seems to bring performance levels nearer to those of BPE models, particularly at greater model scales. This observation can be partly explained by the reduced proportion of total FLOPs consumed by the byte-level Encoder and Decoder, whose scaling with parameter count lags behind that of the Latent Transformer. Specifically, when total parameter count increases by a factor of 20, from 400M to 8B, the local parameters of {\textsc BLT} expand by only about 2x. This is a significant point because enlarging patch sizes impacts the computational cost associated exclusively with the patch Latent Transformer, without affecting the FLOPs attributable to byte-level modules. This phenomenon accounts for why the {\textsc BLT}-Entropy model with patch size 8 increased from 1.6x to 1.7x the Llama 2 parameter count at the larger model scale.

In conclusion, our investigation into patch-length scaling indicates that the {\textsc BLT} patch-based framework attains superior scaling properties when both the patch size and the model capacity are enlarged together. Notably, these favorable scaling dynamics are maintained, and may even be enhanced, as models become larger.

\section{Byte Modeling Improves Robustness} Additionally, we evaluate the robustness of {\textsc BLT} in comparison to token-based models that do not natively process bytes, and we introduce a method for adapting pretrained token-based models to operate at the byte level.
\label{sec:robustness}

\subsection{Character-Level Tasks}

An initial impetus behind the development of byte-level models was their demonstrated resilience to noise at the byte level in input data, as well as their capacity to recognize the compositional structure of tokens—an area where conventional tokenizer-based approaches typically encounter difficulties. To quantitatively assess these capabilities, we conduct supplementary experiments on a set of benchmarks designed to probe both robustness against noisy inputs and sensitivity to character-level information. These evaluations encompass English and multilingual characters, in addition to digits and phonemes. The outcomes of these analyses are summarized in Table~\ref{tab:char_tasks}.

\paragraph{Noisy Data} To assess the resilience of tokenizer-based models in comparison to {\textsc BLT}, we generate perturbed versions of the benchmark classification datasets referenced in Section~\ref{section:evals}. Five different character-level noise schemes are employed to systematically modify the input text: (a) \textit{AntSpeak}: Converts all characters in the text to uppercase and inserts spaces between each character. (b) \textit{Drop}: Eliminates 10\% of characters at random from the text. (c) \textit{RandomCase}: Randomly assigns 50\% of characters to uppercase and 50\% to lowercase throughout the text. (d) \textit{Repeat}: Selects 20\% of the characters and repeats them up to four times. (e) \textit{UpperCase}: Changes all alphabetic characters in the text to uppercase. For evaluation, each noising method is applied separately to either the prompt, the completion, or both, and the resulting performance scores are averaged. Table \ref{tab:char_tasks} summarizes performance on HellaSwag~\citep{Eval_hellaswag} with injected character-level noise, revealing that {\textsc BLT} consistently exceeds the robustness of tokenizer-based approaches, offering an average performance gain of 8 points relative to a comparably trained model, and surpassing the Llama 3.1 variant trained on substantially more data.

\paragraph{Phonology - Grapheme-to-Phoneme (G2P)} We evaluate how effectively {\textsc BLT} can convert sequences of graphemes—letters that make up words—into corresponding phonemic transcriptions that reflect pronunciation. Table \ref{tab:char_tasks} summarizes our findings for the G2P task, conducted in a 5-shot setting with Phonology Bench~\citep{suvarna2024phonologybench}. Our results show that {\textsc BLT} achieves superior performance compared to the baseline Llama 3 1T tokenizer-based approach.

\paragraph{CUTE}
To measure character-level comprehension, we test {\textsc BLT} using the CUTE benchmark~\citep{edman2024cute}, which is composed of a suite of tasks falling into three main groups: composition understanding, recognition of orthographic similarity, and sequence manipulation capabilities. This benchmark proves particularly difficult for many tokenizer-based systems, which often exhibit awareness of token spelling yet face challenges when required to manipulate text on the character level. As shown in Table~\ref{tab:char_tasks}, {\textsc BLT}-Entropy surpasses both BPE Llama 3 models by over 25 points on this suite. Notably, our approach achieves outstanding results on character manipulation, reaching 99.9\% accuracy on both spelling-related tasks. Such substantial gains, despite {\textsc BLT} being trained with only one-sixteenth of the data used for Llama 3.1, suggest that BPE-based models struggle to acquire robust character-level representations. Figure \ref{fig:cute_examples} provides examples highlighting cases where the Llama 3 tokenizer encounters difficulties, whereas {\textsc BLT} handles the tasks effectively. Word deletion and insertion remain the two areas in which BPE models have a performance edge. These forms of word-level manipulation may present greater complexity for byte-level architectures, but the observed performance gap is moderate, implying that constructing word understanding from character representations may be easier than decomposing words into characters. For all experiments, we retain a consistent evaluation protocol, employing original Huggingface prompts. It is possible that the performance of BPE models could improve with tailored prompt engineering.

\paragraph{Low Resource Machine Translation} We assess the performance of {\textsc BLT} on translation tasks involving both directions across six prominent language families and twenty-one low-resource languages, encompassing a variety of writing systems, as represented in the FLORES-101 benchmark~\citep{goyal2022flores}. The corresponding SentencePiece BLEU scores are detailed in Table~\ref{tab:flores}. 

Our findings indicate that {\textsc BLT} surpasses models utilizing the Llama 3 tokenizer, showing an average improvement of 2 BLEU points when translating into English and a 0.5 BLEU point increase when translating from English. For high-resource language pairs, the performance of {\textsc BLT} is generally on par with, or marginally exceeds, that of Llama 3. Notably, in the case of many language pairs from less-resourced language families, {\textsc BLT} achieves superior results relative to Llama 3, highlighting the benefits of byte-level modeling in generalizing across diverse and less frequently seen byte sequences.

\subsection{Training {\textsc BLT} from Llama 3}
\label{sec:distillation}
We investigate an approach wherein {\textsc BLT} models benefit from pre-existing tokenizer-based models to enhance both convergence speed and overall training efficacy. This is accomplished by initializing the global transformer parameters of {\textsc BLT} using those from a pre-trained Llama 3.1 model. In the subsequent training phase, we apply a learning rate for the global transformer that is one-tenth of the rate used for the local encoder and decoder components, following the optimal number of training steps established for Llama 3. A comparative analysis with a baseline {\textsc BLT} model is provided in Table~\ref{tab:distill}. The results clearly demonstrate that initializing {\textsc BLT} with Llama 3.1 weights leads to substantial improvements over both the original Llama 3 and the standard {\textsc BLT} baseline, all trained under equivalent compute budgets. Additionally, when assessed against our {\textsc BLT}-Entropy model in Table~\ref{tab:evals}—which utilized a much larger training set of 1T tokens—{\textsc BLT} initialized from Llama 3.1 still attains higher scores on the MMLU task. These findings indicate that this strategy can provide significant reductions in required training flops while maintaining or improving model performance.

This approach can alternatively be interpreted as adapting models that rely on tokenizers into models that operate without them, in effect transforming a pre-trained LLaMA 3.1 into a {\textsc BLT} architecture. For a thorough assessment, we report results from the original LLaMA 3.1, which was trained on 15T tokens, alongside the corresponding {\textsc BLT} derived from LLaMA 3 in Table \ref{tab:distill}. While our converted model shows only a modest decrease in accuracy on MMLU and HumanEval, it suffers more pronounced reductions on several other benchmarks. These outcomes indicate that additional research is required to maximize the potential of the pre-trained model, especially with respect to optimizing data mixtures and fine-tuning other key hyperparameters.

\section{Ablations and Discussion}
\label{section:ablations}
This section examines ablation experiments that underpin the architectural decisions behind {\textsc BLT}, as well as the selected patching strategy and hyper-parameter configurations used in training the 8B parameter {\textsc BLT} model on the {\textsc BLT}-1T corpus.

\paragraph{Entropy Model Hyper-parameters} To evaluate how the entropy model’s parameter count and context window size impact scaling behavior, we train several byte-level entropy transformers, ranging from 1 million to 100 million parameters, and experiment with context windows between 64 and 512 tokens. Using our $400m$ and $1b$ {\textsc BLT} models trained on the Llama-2 data, we generate bpb versus training flop scaling law plots, which are shown in \autoref{fig:entropy_ablation}. The results reveal that performance improves as both model scale and context window increase, though gains plateau when scaling past 50 million parameters.

\paragraph{Types of Patching}

We conduct an ablation study on the four patching schemes described in Section~\ref{section:patching}: 1) Strided Patching utilizing strides of 4 and 6, 2) Patching at whitespace intervals, 3) BPE Tokenizer patching corresponding to the Llama 3 tokenizer, and 4) Entropy-driven patching leveraging a compact byte-based llm.

Although dynamic patching leads to shorter sequence lengths, we regulate sequence length to ensure that every patching method retains a comparable context length. To avoid potential confounds related to context modeling, every model processes an equivalent expected number of bytes per sequence during both training and inference. The outcomes of these controlled experiments are illustrated in Figure~\ref{fig:scaling-compute-opt}. Among the ablation strategies, all dynamic patching approaches yield better performance than static patching, with space patching demonstrating results nearly matching those of dynamic entropy-based patching.

In \autoref{tab:evals_ablation}, we report benchmark results for {\textsc BLT} models, offering a comparison among tokenizer-based approaches, space patching, and entropy-based patching, each trained for an optimal number of steps on the Llama 2 dataset~\citep{dubey2024llama}. While space patching represents a more straightforward approach that does not require real-time entropy modeling during training, our findings indicate that the improvements associated with entropy-based patching, as noted in the scaling trends of Section~\ref{section:scaling}, also extend to downstream benchmark tasks.\footnote{The space patching results are derived from earlier experiments without cross-attention; however, consistent trends emerge even when cross-attention is incorporated.}

\paragraph{Cross-Attention} 
\autoref{tab:crossattn_ablations_1b} presents an ablation study examining the impact of incorporating cross-attention at different stages within both the encoder and decoder of {\textsc BLT}. For the encoder-side cross-attention, we explore three initialization strategies for the queries: (1) using a shared learned embedding across all global states, (2) utilizing a hash embedding derived from the bytes of each patch, and (3) applying a pooling operation to the encoder’s hidden states corresponding to the patch bytes at the specified encoder layer.

Our results indicate that integrating cross-attention within the \textit{decoder} yields the greatest benefit. In contrast, applying cross-attention in the encoder leads to modest gains, which occur solely when query initialization utilizes pooling. Furthermore, our analysis reveals that cross-attention is especially advantageous for the Common-Crawl dataset, with its effectiveness becoming more pronounced when employing larger patch sizes.

\paragraph{n-gram Hash Embeddings} 

We experiment with n-gram hash embedding vocabulary sizes of 0, 100K, 200K, and 400K, with outcomes summarized in Table~\ref{tab:ngram_ablations_1b}. Our observations reveal that hash embeddings improve performance across all evaluated domains, with pronounced benefits on Wikipedia and Github datasets—showing a 0.04 bpb advantage, as opposed to a 0.01 bpb difference after 15k training steps at the 8B parameter scale. Further, at the 8B level, reducing the number of hashes from 500K to 300K results in only a marginal impact on performance (approximately 0.001 bpb after 15k steps). These findings suggest that hash embeddings play an essential role in enabling {\textsc BLT} to attain comparable results to tokenizer-based approaches; nonetheless, the benefits plateau beyond 300K hashes. Moreover, it is evident that the improvements from hash embeddings and cross-attention are largely independent, each providing enhancements on distinct datasets.

\paragraph{Local Model Hyperparamaters} Table \ref{tab:local_layers} presents an ablation study exploring different configurations for the number of layers in both the local encoder and decoder. With hash n-gram embeddings, {\textsc BLT} demonstrates strong performance when utilizing a very minimal encoder—specifically, a single-layer architecture—combined with a more complex, multi-layer decoder.

\section{Related Work}
\label{section:related_work}

\textbf{Character-Level RNNs:} Character-level language modeling has remained a central focus in neural NLP research since its inception~\citep{sutskever2011generating,mikolov2012subword,graves2013generating}, primarily because these models naturally handle out-of-vocabulary words without relying on explicit back-off techniques. \cite{kim2016character} present an approach where character inputs are processed exclusively through convolutional and highway networks before being passed into LSTM-based RNNs. Their model achieves comparable results with contemporaneous RNN language models for English and demonstrates superior performance on morphologically complex languages, highlighting a notable strength of character-level large language models. Extending these ideas, \cite{kenter2018byte} apply byte-level LSTM architectures for machine comprehension, yielding better results than word-based models, especially for languages like Turkish and Russian that exhibit rich morphological features. In a related study, \cite{zhang2015character} leverage character-level convolutional networks for text classification tasks, showing improved accuracy over word-level approaches in certain scenarios. \citet{chung2019hierarchical} introduce hierarchical LSTM models equipped with boundary-detection mechanisms at various levels, enabling the discovery of latent textual hierarchies and thus further enhancing character-level language model performance. Diverging from attention-based methods, ByteNet by \cite{kalchbrenner2016neural} employs CNN layers operating at the character level for machine translation.

\textbf{Character-Level Transformers:} Transformer architectures leveraging attention mechanisms~\citep{vaswani2017attention} and subword tokenization strategies~\citep{sennrich-etal-2016-neural} have led to notable gains in language modeling and standardized evaluations. Nonetheless, representing text via word or sub-word tokens imposes a specific inductive bias regarding the abstraction level processed by the models. Seeking to unite the strengths of transformer frameworks with earlier encouraging outcomes in character-level language modeling, \citet{al2019character} employ highly deep transformer networks. Augmented by auxiliary loss functions, these models surpass previous LSTM-based character language models. Despite these advances, there remained a marked disparity compared to word-level large language models. Similarly, GPT-2~\citep{radfordgpt2} reported that for expansive corpora such as the 1 Billion Word dataset, language models based on bytes still lagged behind their word-based counterparts.

Although \cite{Choe2019BridgingTG} showed that transformer-based LLMs at the byte level can achieve better performance than subword-based LLMs with a similar number of parameters, such architectures demand substantially higher computational resources and training time. Along related lines, \cite{el2020characterbert} introduced CharFormer, a BERT variant that constructs word-level representations through convolutions over character embeddings. While this method yielded gains, particularly in the medical domain, it also resulted in significant computational overhead. The CANINE model, proposed by \cite{clark2022canine}, is a 150M-parameter encoder-only framework that processes raw character sequences. Its design employs a deep transformer stack—much like the global component in our own work—and incorporates both a local transformer and strided convolution layers for input character downsampling. CANINE surpasses comparable token-level encoder models, such as mBERT, on a range of multilingual downstream tasks. ByT5~\citep{xue2022byt5} investigated byte-level encoder-decoder architectures without any form of patching operations. While ByT5 demonstrated greater resilience to input noise and matched the performance of models that rely on tokenization with one-quarter of the data, the absence of patching necessitated expensive attention computations across all bytes, resulting in extremely high computational costs. The inherent increase in sequence length when modeling at the byte level—as opposed to using subword units—further complicates efficient scaling of such models. More recently, \cite{wang2024mambabyte} leveraged the Mamba Architecture~\citep{gu2023mamba}, which offers fixed-size memory states over extended contexts, to train a byte-level Mamba model without patching. This approach outperformed comparable byte-level transformer models at the 350M parameter level (in bits-per-byte) across various datasets under flop-constrained conditions.

\textbf{Patching-based approaches:} Leveraging patching methods can significantly reduce the computational burden, measured in flops, of byte-level LLMs while potentially maintaining comparable performance. Numerous studies have reported promising results with patching techniques, particularly at smaller scales concerning both model parameters and the volume of training data. For instance, \cite{nawrot-etal-2022-hierarchical} explore static approaches to downsampling and upsampling via patching and introduce the hourglass transformer, which surpasses other byte-level models at the 150M parameter level. Building on this, \cite{nawrot-etal-2023-efficient} advance the methodology with dynamic patching strategies, including an end-to-end trainable boundary predictor, a predictor supervised using particular tokenizers, and an entropy-based model analogous to {\textsc BLT}. Their findings indicate this strategy yields superior performance over standard transformers on language modeling tasks at a scale of 40M parameters with approximately 400M tokens. Further, \cite{lester2024training} analyze training with sequences compressed through arithmetic coding—achieving greater compression than is possible with BPE. By employing an equal-info windows approach, they demonstrate improvements over byte-level baselines for language modeling tasks, though the method still lags behind subword-based alternatives.

Our research is primarily motivated by and closely aligned with MegaByte~\citep{yu2023megabyte}, a decoder-only causal LLM architecture that employs fixed, static patching combined with concatenated representations to transform bytes into patches; a local decoder-side model then reconstructs the byte sequence from these patches. MegaByte has been shown to perform on par with tokenizer-based models when scaled to 1B parameters on datasets of approximately 400B bytes. Throughout our experimental analysis, we include ablations of MegaByte and observe that the static patching approach does not keep pace with the latest compute-efficient tokenizer-based models under equivalent flop budgets. We further illustrate how {\textsc BLT} addresses this performance disparity. Similarly, \cite{slagle2024spacebyte} recognize this limitation with MegaByte and advocate for enhancing the static patching technique by leveraging whitespace and analogous space characters for patching, as well as incorporating a local encoder. Their results reveal that such modifications yield improvements relative to tokenized transformers in compute-matched scenarios across certain data domains, such as Github and arXiv, at the 1B parameter scale. We extend our evaluation to include this model and establish that further advancements in model architecture are necessary for byte-level models to match, and potentially surpass, state-of-the-art token-based systems like Llama 3 at larger scales.

\section{Limitations and Future Work}

In this study, we adopt the number of training steps identified as optimal for Llama~3~\citep{dubey2024llama} when making architectural decisions. It is important to note, however, that these scaling laws were originally established for BPE-level transformer models and might result in less-than-ideal ratios between data and parameter sizes for {\textsc BLT}. Determining scaling laws specific to {\textsc BLT}, which could potentially yield even better scaling properties for our architecture, is left as an avenue for future investigation. Furthermore, a majority of our experiments are limited to models with up to 1B parameters; it remains possible that optimal architectural configurations may shift as we scale to 8B parameters and larger, potentially enhancing performance at greater model sizes.

Current transformer codebases and libraries are primarily optimized for architectures built around tokenizers. Although our study reports experiments matched by theoretical floating point operations and utilizes optimized approaches like FlexAttention to accommodate layers differing from the standard transformer configuration, there remains a possibility that our implementations lag behind tokenizer-based models in terms of real-world execution speed. Consequently, additional improvements could further enhance wall-clock efficiency.

Although {\textsc BLT} employs an entropy model for patching that is trained independently, exploring end-to-end training of the patching mechanism presents a promising avenue for future research. Section \ref{sec:distillation} provides preliminary experimental results suggesting the feasibility of converting large tokenizer-based models, like Llama 3—which are trained on over 10T tokens—into byte-based models. This is achieved by initializing the global transformer with the pretrained weights and keeping them fixed. Continued investigation in this direction could reveal strategies that not only maintain the strengths of byte-based models, but also surpass the capabilities of tokenizer-based counterparts, all without necessitating training from the ground up.

\section{Conclusion}
\label{section:conclusion}

In this work, we introduce the Byte Latent Transformer (\textbf{BLT}), an innovative model architecture that challenges the traditional reliance on predetermined token vocabularies in large language models. {\textsc BLT} leverages a flexible, trainable approach to aggregate bytes into variable-length patches, thereby distributing computation in accordance with the underlying data’s intricacy. This strategy yields notable gains in computational efficiency and model robustness. Through an extensive scaling analysis, we show that {\textsc BLT} models are able to achieve performance on par with established tokenization-centric models such as Llama 3 across configurations up to 8B parameters and 4T bytes. Additionally, {\textsc BLT} is capable of reducing inference flops by as much as 50\%, with only modest impacts on evaluation performance. A key advantage of {\textsc BLT} is its ability to simultaneously expand model and patch sizes under a constant inference computational budget, offering a new scaling dimension that is well-suited to the computational environments prevalent in real-world applications. By operating directly on raw bytes, {\textsc BLT} not only bolsters resilience against atypical or noisy inputs but also enhances the modeling of sub-word units. Collectively, these advances establish {\textsc BLT} as a compelling and flexible alternative to tokenization-based language models, enabling greater adaptability, scalability, and robustness.

\end{document}